\documentclass[10pt,letterpaper]{article}
\usepackage[T1]{fontenc}
\usepackage[utf8]{inputenc}
\usepackage{lmodern}
\usepackage[margin=0.76in]{geometry}
\usepackage[table]{xcolor}
\usepackage{array,tabularx,longtable,booktabs}
\usepackage{enumitem}
\usepackage{amssymb,amsmath}
\usepackage{hyperref}
\usepackage{fancyhdr,lastpage}
\usepackage{microtype}
\usepackage{parskip}
\definecolor{cdblue}{HTML}{203957}
\definecolor{cdgray}{HTML}{586677}
\hypersetup{colorlinks=true,linkcolor=cdblue,urlcolor=cdblue,pdftitle={Chapter 01: The Problem of Delivering Software}}
\pagestyle{fancy}
\fancyhf{}
\lhead{\small\color{cdgray}Continuous Delivery Engineering}
\rhead{\small\color{cdgray}Study guide}
\cfoot{\thepage\ / \pageref{LastPage}}
\setlength{\headheight}{14pt}
\setlength{\emergencystretch}{2em}
\setlist[itemize]{leftmargin=*,itemsep=3pt,topsep=4pt}
\setlist[enumerate]{leftmargin=*,itemsep=3pt,topsep=4pt}
\begin{document}
\begin{center}
{\LARGE\bfseries\color{cdblue}Chapter 01: The Problem of Delivering Software}\\[7pt]
{\large I: Foundations}\\[6pt]
{\small Reference edition: 2026-10-30 $\bullet$ Original educational and operational material}
\end{center}
\noindent\color{cdblue}\rule{\linewidth}{1.2pt}\color{black}
\vspace{0.4em}
\section{Chapter scope and provenance}
Primary reading: Humble and Farley, Continuous Delivery (2010), Chapter 1, "The Problem of Delivering Software"; Part I: Foundations. This guide is an original synthesis and not a reproduction of the chapter. The book provides the foundational delivery principles; current AWS, Terraform and GitHub Actions references in the applied sections are supplementary.


\section{Learning outcomes}
\begin{enumerate}
\item Identify manual-deployment and late-integration risk.
\item Explain why each change should produce a potential release candidate.
\item Quantify lead time and time-to-feedback.
\end{enumerate}

\section{Core conceptual model}
Focus areas: release anti-patterns; automation and feedback; releasability as a shared responsibility.

\begin{itemize}
\item Input: a version-controlled change, explicit requirements, and a reproducible execution context.
\item Mechanism: controlled feedback, automation, evidence capture, and a stop condition where risk exceeds tolerance.
\item Output: a verifiable release candidate or a diagnosed failure with an accountable owner.
\end{itemize}

\section{Reading and analysis protocol}
\begin{enumerate}
\item Read the chapter and record five propositions, distinguishing an author claim from an example or assumption.
\item Identify two failure modes, one recovery technique, and one measurable acceptance condition.
\item Trace each technique to the stages and controls in the deployment-pipeline-architecture reference.
\item State how the proposed practice changes for a regulated system with segregated duties.
\end{enumerate}

\section{Detailed technical study notes}
\subsection{Delivery bottlenecks}
Model the release process as a flow of work and evidence. Queue time and correction work often dominate the active build time; measuring only build duration hides the constraint.

\subsection{Manual release failure modes}
Unscripted instructions introduce variation, unverifiable state transitions, unrecorded expertise and inconsistent recovery paths. Demonstrate which manual step can be made repeatable.

\subsection{Production-like validation}
Late discovery of environment differences shifts integration risk to the most expensive point in the cycle. Decide how much parity is necessary for each dependency.

\subsection{Rapid feedback}
Feedback has no economic value until a responsible engineer can act on it. Set a target for time-to-first-actionable-failure and define the response owner.

\subsection{Release candidate semantics}
A candidate is an identified artifact and evidence set, not simply a Git branch name. Promotion should not change its identity.

\subsection{Organization-wide responsibility}
Developers, testers, database engineers, security and operators each own parts of releasability; avoid assigning all release risk to a final operator.

\section{Domain-specific challenge questions}
\begin{enumerate}
\item For delivery bottlenecks, identify the precondition, expected evidence, adversarial failure case and measurable corrective action.
\item For manual release failure modes, identify the precondition, expected evidence, adversarial failure case and measurable corrective action.
\item For production-like validation, identify the precondition, expected evidence, adversarial failure case and measurable corrective action.
\item For rapid feedback, identify the precondition, expected evidence, adversarial failure case and measurable corrective action.
\end{enumerate}
\section{Guided exercise}
Draw a value-stream map from commit to production; identify three waits and two manual handoffs.

\subsection{Required evidence}
\begin{itemize}
\item A version-controlled design or executable artifact.
\item A failure demonstration, negative test or counterexample.
\item An improvement statement with baseline, target, owner and verification method.
\end{itemize}


\section{Engineering scenario}
A release takes 18 hours, needs four specialists, and lacks recovery rehearsal. Design a low-risk first intervention.

\begin{itemize}
\item Describe the failure propagation path and the earliest feasible detection point.
\item Recommend a minimal remediation and its rollback or containment plan.
\item Define a measure that would demonstrate risk reduction.
\end{itemize}

\section{Knowledge checks}
\begin{enumerate}
\item Which condition is the gate, and which action is the response when it fails?
\item What evidence distinguishes successful automation from an assumed successful release?
\item Where might environments, data, dependencies or people invalidate the technique?
\item What feedback could be moved earlier without losing necessary confidence?
\item Which assumptions are book principles, and which are modern supplementary implementation choices?
\end{enumerate}

\section{Completion rubric}
Score each dimension from 0 (missing) to 4 (demonstrated and independently verifiable): conceptual accuracy; reproducibility; negative testing; security and safety; documentation and traceability. Pass threshold: 16/20 with no zero on reproducibility or safety.


\section{Prerequisites and onward route}
Git, basic tests, and a build runner. Next: Chapters 2 onward.



\end{document}
